\documentclass[10pt,a4paper]{amsart}
\usepackage{amsmath,amssymb,amsthm,graphicx}
\usepackage{verbatim}
\newenvironment{reference}{\comment}{\endcomment}
%\newenvironment{reference}{}{}
\newenvironment{slogan}{\comment}{\endcomment}
\newenvironment{history}{\comment}{\endcomment}
% For commutative diagrams we use Xy-pic
\usepackage[all]{xy}

% We use 2cell for 2-commutative diagrams.
\xyoption{2cell}
% Theorem environments.
%
\theoremstyle{plain}
\newtheorem{theorem}[subsection]{Theorem}
\newtheorem{proposition}[subsection]{Proposition}
\newtheorem{lemma}[subsection]{Lemma}

\theoremstyle{definition}
\newtheorem{definition}[subsection]{Definition}
\newtheorem{example}[subsection]{Example}
\newtheorem{exercise}[subsection]{Exercise}
\newtheorem{situation}[subsection]{Situation}

\theoremstyle{remark}
\newtheorem{remark}[subsection]{Remark}
\newtheorem{remarks}[subsection]{Remarks}

\numberwithin{equation}{subsection}

% Macros
%
\def\lim{\mathop{\mathrm{lim}}\nolimits}
\def\colim{\mathop{\mathrm{colim}}\nolimits}
\def\Spec{\mathop{\mathrm{Spec}}}
\def\Hom{\mathop{\mathrm{Hom}}\nolimits}
\def\Ext{\mathop{\mathrm{Ext}}\nolimits}
\def\SheafHom{\mathop{\mathcal{H}\!\mathit{om}}\nolimits}
\def\SheafExt{\mathop{\mathcal{E}\!\mathit{xt}}\nolimits}
\def\Sch{\mathit{Sch}}
\def\Mor{\mathop{\mathrm{Mor}}\nolimits}
\def\Ob{\mathop{\mathrm{Ob}}\nolimits}
\def\Sh{\mathop{\mathit{Sh}}\nolimits}
\def\NL{\mathop{N\!L}\nolimits}
\def\CH{\mathop{\mathrm{CH}}\nolimits}
\def\proetale{{pro\text{-}\acute{e}tale}}
\def\etale{{\acute{e}tale}}
\def\QCoh{\mathit{QCoh}}
\def\Ker{\mathop{\mathrm{Ker}}}
\def\Im{\mathop{\mathrm{Im}}}
\def\Coker{\mathop{\mathrm{Coker}}}
\def\Coim{\mathop{\mathrm{Coim}}}

% Boxtimes
%
\DeclareMathSymbol{\boxtimes}{\mathbin}{AMSa}{"02}

%
% Macros for moduli stacks/spaces
%
\def\QCohstack{\mathcal{QC}\!\mathit{oh}}
\def\Cohstack{\mathcal{C}\!\mathit{oh}}
\def\Spacesstack{\mathcal{S}\!\mathit{paces}}
\def\Quotfunctor{\mathrm{Quot}}
\def\Hilbfunctor{\mathrm{Hilb}}
\def\Curvesstack{\mathcal{C}\!\mathit{urves}}
\def\Polarizedstack{\mathcal{P}\!\mathit{olarized}}
\def\Complexesstack{\mathcal{C}\!\mathit{omplexes}}
% \Pic is the operator that assigns to X its picard group, usage \Pic(X)
% \Picardstack_{X/B} denotes the Picard stack of X over B
% \Picardfunctor_{X/B} denotes the Picard functor of X over B
\def\Pic{\mathop{\mathrm{Pic}}\nolimits}
\def\Picardstack{\mathcal{P}\!\mathit{ic}}
\def\Picardfunctor{\mathrm{Pic}}
\def\Deformationcategory{\mathcal{D}\!\mathit{ef}}
\usepackage{cleveref}
\crefname{originalreference}{original reference}{original references}
\Crefname{originalreference}{Original reference}{Original references}
\makeatletter
\newlabel{algebraic-lemma-characterize-representable-by-algebraic-spaces}{{algebraic-lemma-characterize-representable-by-algebraic-spaces}{source}{}{sourcepointer-1}{}}
\newlabel{algebraic-lemma-characterize-representable-by-algebraic-spaces@cref}{{[originalreference][0][]algebraic-lemma-characterize-representable-by-algebraic-spaces}{[1][0][]source}{}{}{}}
\makeatother
\crefname{originallemma}{original lemma}{original lemmas}
\Crefname{originallemma}{Original lemma}{Original lemmas}
\makeatletter
\newlabel{lemma-inertia}{{lemma-inertia}{source}{}{sourcepointer-2}{}}
\newlabel{lemma-inertia@cref}{{[originallemma][0][]lemma-inertia}{[1][0][]source}{}{}{}}
\makeatother
\title{Representability by algebraic spaces is equivalent to the second diagonal being an isomorphism}
\author{The Stacks Project Authors}
\date{stacks-project:stacks-morphisms.tex:local-snapshot-20261012; GNU FDL 1.2 or later}
\begin{document}
\maketitle
\noindent\textit{Editorial note.} All delivered mathematical text is the author's own, copied from the pinned Stacks Project chapter source sources/stacks-algebraic-stacks-morphisms/pinned-original-source.tex; no statement or proof step was rewritten, reordered or removed.\par\medskip
\noindent\textit{Editorial note.} Preamble provenance: the Stacks Project keeps one single shared preamble for every chapter; this card takes it directly from the staged sources/stacks-algebraic-stacks-morphisms/preamble.tex as three verbatim spans (lines 8-12, 14-18 and 154-217), so no derived file was written and nothing was added or redefined. The wrapper class and the base packages are supplied by the assembler.\par\medskip
\noindent\textit{Editorial note.} Bibliography provenance: neither the counted statement nor its proof contains a citation (this lane rescans the whole span before writing the card), so the card carries no bibliography and no entry was reconstructed.\par\medskip
\noindent\textit{Editorial note.} Selection provenance: this card was opened under the file-based control contract (CONTROL.md), with the mechanical pre-card gate run on its statement and proof interval before the card was written.\par\medskip
\noindent\textit{Allowed standard prerequisites (editorial).}\begin{itemize}
\item The characterisation of morphisms representable by algebraic spaces through faithfulness, and the description of the relative inertia as the kernel of the functor on isomorphism sheaves. Invoked inside the counted proof (source lines 1717-1734) as pointers to results of the chapter on algebraic stacks. Translates representability into a condition on the kernel of the map on Isom sheaves and identifies that kernel with a fibre product of the inertia.
\item The description of the second diagonal as the identity section of the relative inertia group stack. The statement uses the second diagonal and the relative inertia interchangeably, and the proof invokes the lemma identifying the kernel of the Isom map with a fibre product of the inertia. Provides the object whose triviality is equivalent to the conditions of the statement.
\end{itemize}\medskip
\ifdefined\hypertarget\hypertarget{sourcepointer-1}{}\fi
\noindent\textit{Source reference.} Original reference algebraic-lemma-characterize-representable-by-algebraic-spaces: sources/stacks-algebraic-stacks-morphisms/pinned-original-source.tex, source lines 1720--1720. Navigation pointer only: the counted text refers to the Stacks item algebraic-lemma-characterize-representable-by-algebraic-spaces, which belongs to another chapter of the project and is not part of this card. The pointer carries the target's own label as its number because this chapter's staged snapshot does not contain the other chapter's tag table; its content is not reproduced here.\par\medskip
\ifdefined\hypertarget\hypertarget{sourcepointer-2}{}\fi
\noindent\textit{Source reference.} Original lemma lemma-inertia: sources/stacks-algebraic-stacks-morphisms/pinned-original-source.tex, source lines 1327--1327. Navigation pointer only: the counted text refers to the Stacks item lemma-inertia, which belongs to another chapter of the project and is not part of this card. The pointer carries the target's own label as its number because this chapter's staged snapshot does not contain the other chapter's tag table; its content is not reproduced here.\par\medskip
\begin{lemma}
\label{lemma-second-diagonal}
Let $f : \mathcal{X} \to \mathcal{Y}$ be a morphism of algebraic stacks.
The following are equivalent:
\begin{enumerate}
\item the morphism $f$ is representable by algebraic spaces,
\item the second diagonal of $f$ is an isomorphism,
\item the group stack $ \mathcal{I}_{\mathcal{X}/\mathcal{Y}}$
is trivial over $\mathcal X$, and
\item for a scheme $T$ and a morphism $x : T \to \mathcal{X}$
the kernel of $\mathit{Isom}_\mathcal{X}(x, x) \to
\mathit{Isom}_\mathcal{Y}(f(x), f(x))$ is trivial.
\end{enumerate}
\end{lemma}

\begin{proof}
We first prove the equivalence of (1) and (2).
Namely, $f$ is representable by algebraic spaces if and only if $f$ is
faithful, see
Algebraic Stacks,
Lemma \ref{algebraic-lemma-characterize-representable-by-algebraic-spaces}.
On the other hand, $f$ is faithful if and only if for every object $x$
of $\mathcal{X}$ over a scheme $T$ the functor $f$ induces an injection
$\mathit{Isom}_\mathcal{X}(x, x) \to
\mathit{Isom}_\mathcal{Y}(f(x), f(x))$,
which happens if and only if the kernel $K$ is trivial, which happens if and
only if $e : T \to K$ is an isomorphism for every $x : T \to \mathcal{X}$.
Since $K = T \times_{x, \mathcal{X}} \mathcal{I}_{\mathcal{X}/\mathcal{Y}}$
as discussed above, this proves the equivalence of (1) and (2). To prove
the equivalence of (2) and (3), by the discussion above, it suffices to
note that a group stack is trivial if and only if its identity section
is an isomorphism. Finally, the equivalence of (3) and (4) follows
from the definitions: in the proof of Lemma \ref{lemma-inertia}
we have seen that the kernel in (4) corresponds to the fibre product
$T \times_{x, \mathcal{X}} \mathcal{I}_{\mathcal{X}/\mathcal{Y}}$ over $T$.
\end{proof}

\end{document}
